\clearpage
\section{Środowisko eksperymentalne i odtwarzalność}
\label{chap:srodowisko}

Badanie główne obejmuje nominalnie 3132 zadania konfiguracji i~14\,976 dopasowań sieci uczonych w~foldach. Przy budżecie godziny na konfigurację nawet pojedyncze niepowodzenie może oznaczać istotny koszt powtórzenia obliczeń. Środowisko przygotowano więc tak, aby ukończone wyniki można było zachować przy wznowieniu kampanii, a~raporty odtworzyć z~zapisanych predykcji. Każdy wynik pozostaje powiązany z~danymi, ustawieniami i~warunkami wykonania.

\subsection{Wymagania wynikające z metodyki}
\label{sec:wymagania-platformy}

Platforma musi zapewniać spójne mapowanie danych i~klas, kontrolę podziałów, odtwarzalne generowanie konfiguracji oraz wybór modeli bez korzystania z~etykiet testowych. Powinna ponadto umożliwiać przerwanie długiej kampanii, rozpoznanie kompletnych wyników i~odtworzenie raportów bez ponownego uczenia. Te wymagania wynikają z~liczby konfiguracji i~z~konieczności przypisania każdej obserwacji naukowej do konkretnego wykonania.

Przyjęto lokalne, sekwencyjne uruchamianie konfiguracji i~foldów w~jednym procesie sterującym. Ustawienia z~plików YAML są odczytywane i~sprawdzane przez centralny moduł konfiguracji. Zachowano obiektowe adaptery modeli oraz wspólny kontrakt predykcji prawdopodobieństw. Umożliwia to współistnienie lokalnych implementacji, natywnych modeli bibliotecznych i~stałego modelu pretrenowanego.

Wymaganiem naukowym jest również jawność niepowodzeń. Awaria nie może zostać przekształcona w~nieoznaczony brak lub pozornie niski koszt. Zapis próby obejmuje stan, czas i~diagnostykę, natomiast reguły raportowania odróżniają jawnie nieudaną próbę od nieukończonego zadania.

\subsection{Architektura i integracje}
\label{sec:architektura-platformy}

Za uruchamianie modeli i~walidację odpowiadają biblioteki TabArena i~AutoGluon w~ustalonych wersjach. Własny kod łączy je z~konfiguracją badania, obsługą danych oraz zapisem wyników. Zakres tej integracji i~odpowiedzialność poszczególnych modułów przedstawia tabela~\ref{tab:warstwy-platformy}.

\begin{table}[tbp]
\centering\small
\caption{Warstwy platformy i~ich odpowiedzialność. Nazwy modułów wskazują źródła implementacji, a~nie niezależne usługi.}
\label{tab:warstwy-platformy}
\begin{tabularx}{\linewidth}{P{0.22\linewidth}P{0.27\linewidth}Y}
\toprule
Warstwa & Główne moduły & Odpowiedzialność \\
\midrule
Konfiguracja & \path{configuration.py}, \path{configs/} & Typy, walidacja i~łączenie ustawień z~plików oraz nadpisań \\
Dane & \path{data/tasks.py} & Przygotowanie, manifesty, mapowanie klas i~zewnętrzne podziały \\
Katalog i~przestrzenie & \path{catalog.py}, \path{specification.py} & Możliwości metod i~losowanie konfiguracji \\
Wykonanie & \path{backend.py}, adaptery modeli & Uczenie foldów, OOF i~prawdopodobieństwa testowe \\
Zasoby & \path{runtime.py} & Walidacja urządzenia, wątki i~kontekst dopasowania \\
Zapis i~wznawianie & \path{runner.py}, \path{artifacts.py} & Tożsamość, blokady, próby, integralność i~wznowienie \\
Raportowanie & \path{reporting.py}, \path{evaluation/metrics.py} & Wybór OOF, zespoły oraz raporty JSON i~CSV \\
\bottomrule
\end{tabularx}
\end{table}

Główna ścieżka wykorzystuje sześć lokalnych integracji: MLP, TabNet, FT-Transformer, regresję logistyczną, XGBoost i~CatBoost. KNN, RealMLP, ModernNCA oraz TabM są uruchamiane przez natywne ścieżki przypiętych bibliotek. Osobny adapter obsługuje stały TabPFN. Wspólny katalog eksperymentów łączy te ścieżki, dzięki czemu dodanie natywnej integracji nie wymaga powielania jej procedury uczenia w~lokalnym adapterze.

Moduł wykonawczy otrzymuje dane z~zachowaniem ról cech numerycznych i~kategorycznych. Automatyczne wspólne przygotowanie cech jest wyłączone, a~za przekształcenia odpowiada właściwa integracja. Przy GPU pomiar dopasowania i~predykcji jest otoczony synchronizacją urządzenia. Ogranicza to błąd wynikający z~asynchronicznego wykonania operacji, ale nie usuwa konieczności precyzyjnego określenia granic pomiaru.

Rysunek~\ref{fig:architektura} przedstawia zależności komponentów i~oddzielenie uczenia od wtórnego odczytu zapisanych predykcji. Interfejs wiersza poleceń znajduje się w~\path{experiments/cli.py}; konfiguracja jest rozwiązywana przed wywołaniem właściwego polecenia.

\begin{figure}[tbp]\centering\small
% Opracowanie własne; schemat procedury opisanej w rozdziałach 3–5.
\begin{minipage}[t]{0.31\linewidth}\centering\textbf{Przygotowanie}\par\medskip
\fbox{\parbox{0.85\linewidth}{\centering Konfiguracja YAML}}\par\smallskip$\downarrow$\par\smallskip\fbox{\parbox{0.85\linewidth}{\centering Dane i~podziały}}\par\smallskip$\downarrow$\par\smallskip\fbox{\parbox{0.85\linewidth}{\centering Katalog konfiguracji}}\end{minipage}
\hfill
\begin{minipage}[t]{0.31\linewidth}\centering\textbf{Wykonanie}\par\medskip
\fbox{\parbox{0.85\linewidth}{\centering TabArena i~AutoGluon}}\par\smallskip$\downarrow$\par\smallskip\fbox{\parbox{0.85\linewidth}{\centering Adapter wybranej metody}}\par\smallskip$\downarrow$\par\smallskip\fbox{\parbox{0.85\linewidth}{\centering Uczenie i~predykcje}}\end{minipage}
\hfill
\begin{minipage}[t]{0.31\linewidth}\centering\textbf{Trwały zapis}\par\medskip
\fbox{\parbox{0.85\linewidth}{\centering Manifest i~historia prób}}\par\smallskip$\downarrow$\par\smallskip\fbox{\parbox{0.85\linewidth}{\centering Predykcje i~metadane}}\par\smallskip$\downarrow$\par\smallskip\fbox{\parbox{0.85\linewidth}{\centering Raport D/T/T+E}}\end{minipage}\par\medskip Przepływ: przygotowanie $\longrightarrow$ wykonanie $\longrightarrow$ zapis.\par Wznowienie odczytuje ukończone zadania; raportowanie korzysta z~zapisanych predykcji.

\caption{Schemat lokalnej platformy oraz przepływu artefaktów; opracowanie własne.}
\label{fig:architektura}
\end{figure}

\subsection{Izolacja informacji i identyfikowalność danych}
\label{sec:izolacja-platformy}

Przygotowane dane są przechowywane jako plik Parquet z~manifestem i~sumami kontrolnymi. Źródła OpenML mają określone identyfikatory i~wersje; tożsamość pobieranych archiwów UCI określa suma kontrolna zawartości. Santander jest pobierany automatycznie z~OpenML (identyfikator 46859, wersja~2) w~ramach wspólnej procedury przygotowania danych. Zgodność schematu i~hashu chroni przed niezamierzoną podmianą danych, lecz nie rozstrzyga problemów populacji, przecieku lub duplikatów wskazanych w~sekcji~\ref{sec:dane-audyt}.

Zewnętrzne podziały zachowują identyfikatory wierszy, etykiety i~grupy pacjentów. Kolumny grupujące nie wchodzą do predyktorów. Backend zapisuje rzeczywiste indeksy walidacji wewnętrznej, co umożliwia późniejszą kontrolę pokrycia OOF i~rozłączności pacjentów. Zdefiniowane są także zastrzeżone nazwy pomocniczych kolumn; konflikt z~danymi wejściowymi jest odrzucany zamiast prowadzić do cichego nadpisania cechy.

Łączenie predykcji wymaga zgodności zarówno obserwacji, jak i~znaczenia kolumn. Macierze o~tych samych wymiarach mogą bowiem opisywać inne wiersze lub mieć odwróconą kolejność klas. Mapowanie klas pochodzi z~definicji przygotowanego zadania i~jest zapisywane wraz z~wynikiem. Podczas zapisu sprawdza się zgodność etykiet z~identyfikatorami wierszy, wymiary macierzy oraz skończoność, zakres i~sumy prawdopodobieństw. Niewielkie odchylenia numeryczne mogą być normalizowane. Przed zbudowaniem zespołu porównuje się ponadto etykiety i~identyfikatory wierszy OOF oraz testu. Raporter nie porównuje osobno list klas między rekordami, dlatego poprawność ich kolejności opiera się na wspólnym mapowaniu zadania; sama kontrola numeryczna jej nie potwierdza.

Ustawienie pola czasowego w~definicji zadania prowadzi do jawnego odrzucenia niewspieranego podziału czasowego. Kontrola techniczna blokuje zatem niewspieraną procedurę, ale nie zastępuje naukowego uzasadnienia podziału IID.

Raport zachowuje odziedziczoną nazwę pola \texttt{log\_loss\_val}, chociaż zawarta w~nim wartość jest stratą na zewnętrznym teście. Przy przenoszeniu wyników do tabel pracy rozstrzygająca jest populacja powstania predykcji. Podobnie zbiorcze binarne AP odnosi się do klasy zakodowanej jako~1; do oceny rzadkiej klasy służy AP zapisane osobno dla każdej klasy w~\texttt{per\_class}. Kolumnę wybiera się na podstawie mapowania rzeczywistej etykiety zdarzenia, a~nie założenia, że ma ona kod~1. Te rozbieżności między nazwą techniczną a~interpretacją pomiaru muszą być uwzględnione w~analizie.

\subsection{Artefakty, awarie i wznowienie}
\label{sec:artefakty-wznowienie}

Artefakty zestawione w~tabeli~\ref{tab:artefakty} pozwalają wznowić obliczenia bez pomieszania wyników różnych konfiguracji. Identyfikator badania jest wyznaczany z~efektywnych ustawień, ścieżek, manifestów danych, skrótu źródeł i~pliku zależności, wybranych wersji pakietów oraz zapisanej charakterystyki sprzętu. Dane są serializowane w~ustalonej postaci i~identyfikowane przez SHA-256. Skrót źródeł obejmuje moduły wykonawcze, \path{pyproject.toml} i~\path{uv.lock}; nie obejmuje automatycznie dokumentacji, testów ani skryptów wdrożeniowych. Identyfikator commitu i~obrazu kontenera należy zachować osobno. Zmiana ścieżek może zmienić tożsamość mimo tej samej treści danych.

\begin{table}[tbp]
\centering\small
\caption{Podstawowe artefakty jednego badania. Ścieżki odnoszą się do katalogu jego tożsamości.}
\label{tab:artefakty}
\begin{tabularx}{\linewidth}{P{0.46\linewidth}Y}
\toprule
Plik lub katalog & Znaczenie \\
\midrule
\path{manifest.json}, \path{index.json} & Zamrożone ustawienia i~wykaz zadań \\
\path{splits/DATASET.json} & Tożsamość zewnętrznego podziału \\
\path{jobs/DATASET-MODEL-000/manifest.json} & Konkretna konfiguracja \\
\path{jobs/.../status.json}, \path{attempts/UUID.json} & Stan i~historia prób \\
\path{jobs/.../records/result.json} & Metadane wykonania oraz czasy \\
\path{jobs/.../predictions/result.npz} & Wyrównane predykcje OOF i~testowe; TabPFN bez OOF \\
\path{failure.json}, \path{failures.json} & Diagnostyka niepowodzeń we właściwym zakresie \\
\path{report.json}, \path{report.csv} & Pochodne wyniki D, T i~T+E \\
\bottomrule
\end{tabularx}
\end{table}

Sama konfiguracja nie wystarcza do odtworzenia warunków uczenia, ponieważ rozmiar partii zależy również od danych danego foldu i~ograniczeń wykonania. Dlatego dla lokalnych MLP, FT-Transformera i~TabNet w~bloku \texttt{hyperparameters\_fit} zapisuje się efektywny rozmiar partii (\texttt{effective\_batch\_size}), a~dla TabNet także skorygowany rozmiar partii wirtualnej (\texttt{effective\_virtual\_batch\_size}). Wartości te pozwalają ustalić, z~jakimi partiami rzeczywiście pracował model po zastosowaniu reguł z~sekcji~\ref{sec:strojenie}.

Ten sam blok zawiera znacznik zakończenia z~powodu budżetu (\texttt{budget\_truncated}), liczbę ukończonych epok (\texttt{completed\_epochs}), przydzielony czas (\texttt{granted\_seconds}) i~czas dopasowania mierzony przez lokalną integrację (\texttt{fit\_seconds\_wrapper}). Zapis obejmuje więc zarówno ustawienia wykonania, jak i~przebieg uczenia. Czas lokalnego dopasowania dotyczy pojedynczego foldu; koszt całej konfiguracji i~wyszukiwania rozlicza się zgodnie z~sekcją~\ref{sec:koszt-metodyka}.

Metadane JSON i~predykcje są najpierw zapisywane do plików tymczasowych. Po opróżnieniu buforów zapisu i~synchronizacji z~dyskiem plik tymczasowy atomowo zastępuje plik docelowy. Blokada wyłączna zapobiega jednoczesnemu zapisowi tego samego katalogu przez dwa uruchomienia. Przed pominięciem ukończonego zadania moduł wykonawczy kontroluje spójność statusu, rekordów, predykcji i~ich sum kontrolnych. Sam fakt istnienia pliku nie wystarcza do uznania zadania za ukończone.

Awarie są zapisywane wraz z~diagnostyką, a~niezależne zadania mogą być kontynuowane. Wystąpienie niepowodzeń powoduje niezerowy kod zakończenia głównego uruchomienia. Automatyczny raport powstaje po jego pomyślnym ukończeniu. Oddzielne polecenie raportowania może obsłużyć jawnie nieudane konfiguracje losowe, gdy wszystkie zadania zostały podjęte; nieudany wariant domyślny albo zadanie nieukończone uniemożliwiają pełny raport.

Po próbie usuwane są dopasowane modele, a~pozostają predykcje i~metadane. Umożliwia to ponowne wyliczenie raportu bez zachowania wszystkich sieci. Nie umożliwia jednak wznowienia w~środku foldu ani wykorzystania usuniętego modelu do nowych predykcji. Konfigurację przerwaną bez kompletnego zapisu wyników trzeba dopasować ponownie w~całości; poprawne zakończenie uczenia wskutek ograniczenia budżetu nie wymaga takiego powtórzenia. Tożsamość badania i~kontrole zgodności ograniczają również automatyczne przenoszenie wyników do zmienionego eksperymentu.

\subsection{Dobór partii i ograniczenia pamięci}
\label{sec:partie-platformy}

W~MLP, FT-Transformerze i~TabNet rozmiar partii wynika z~liczności treningowej części danego foldu. Strojenie obejmuje mnożnik tej bazowej wartości, wybierany ze zbioru $\{0{,}5;\,1;\,2\}$; w~konfiguracji domyślnej wynosi on~1. Dla MLP i~FT-Transformera podstawą doboru jest liczba przykładów treningowych~$n$. TabNet uwzględnia dodatkowo kategorie: progi stosuje się do wartości $n'=n(1+C/5000)$, gdzie $C$ oznacza sumę liczności słowników kategorii utworzonych na treningu danego foldu, wraz z~pozycjami zarezerwowanymi dla nieznanych wartości. Przyjęte progi zestawiono w~tabeli~\ref{tab:dobor-partii}. Jest to heurystyka wykonania, której nie należy utożsamiać z~empirycznie potwierdzonym optymalnym doborem partii.

\begin{table}[tbp]
\centering\small
\caption{Bazowe rozmiary partii w~lokalnych integracjach. Próg odnosi się do $n$ dla MLP i~FT-Transformera oraz do $n'$ dla TabNet.}
\label{tab:dobor-partii}
\begin{tabular}{lrrr}
\toprule
Przedział & MLP i~FT-Transformer & TabNet & Partia wirtualna TabNet \\
\midrule
$[0, 2000)$ & 64 & 128 & 128 \\
$[2000, 8000)$ & 128 & 256 & 128 \\
$[8000, 32000)$ & 256 & 1024 & 128 \\
$[32000, 128000)$ & 512 & 2048 & 256 \\
$[128000, \infty)$ & 1024 & 4096 & 512 \\
\bottomrule
\end{tabular}
\end{table}

Po zastosowaniu mnożnika rozmiar partii zaokrągla się w~dół do potęgi dwójki, z~minimum ośmiu przykładów. Następnie uwzględnia się ograniczenie pamięci FT-Transformera opisane w~sekcji~\ref{sec:koszt-metodyka} lub korektę TabNet zapobiegającą pozostawieniu jednoelementowej partii. W~TabNet osobno dobierana jest partia wirtualna, czyli rozmiar fragmentów partii wykorzystywanych przy normalizacji. Jej wartość bazowa pochodzi z~tej samej tabeli, lecz nie jest skalowana mnożnikiem. Ogranicza się ją do rozmiaru partii uczenia, po czym w~razie potrzeby zwiększa, aby uniknąć jednoelementowych fragmentów także w~ostatniej, niepełnej partii. Jawne zadanie całkowitego rozmiaru partii wyłącza automatyczny dobór i~działanie mnożnika, ale pozostawia wymienione ograniczenia wykonania.

W~FT-Transformerze efektywny rozmiar partii jest dodatkowo ograniczany na podstawie liczby cech i~głów uwagi. Dla $d$ predyktorów i~$H$ głów przyjmuje się $4H(d+1)^2$ bajtów na przykład dla pojedynczego tensora wyników uwagi w~formacie float32, z~uwzględnieniem tokenu klasyfikacyjnego. Z~przydziału 4~GiB wyznacza się limit liczby przykładów, zaokrąglany w~dół do potęgi dwójki, z~dolnym ograniczeniem tego limitu do ośmiu. Efektywna partia jest mniejszą z~wartości uzyskanej po automatycznym doborze i~zastosowaniu mnożnika oraz tak wyznaczonego limitu. Jeżeli rozmiar zadano bezpośrednio jako liczbę całkowitą, limit nakłada się na tę liczbę. Przy 1776 predyktorach limit wynosi 32 przykłady dla ośmiu głów oraz 64 dla czterech. Ograniczenie obowiązuje także podczas walidacji i~predykcji. Jest to heurystyka pamięciowa dotycząca jednego tensora, nie gwarancja zmieszczenia pełnego przebiegu w~VRAM. Ponieważ zmniejszenie partii zmienia również liczbę aktualizacji parametrów na epokę, efektywny rozmiar należy uwzględniać przy interpretacji zarówno kosztu, jak i~jakości modelu.

\subsection{Pomiar kosztu i interpretacja raportu}
\label{sec:pomiary-platformy}

Koszt wyszukiwania i~suma czasów dopasowania konfiguracji wybranych do użycia opisują różne etapy procedury. Czas \texttt{search\_fit\_seconds} obejmuje udane dopasowania, natomiast \texttt{search\_wall\_seconds} uwzględnia zapisane próby, także zakończone błędem i~powtórzenia. Przerwanie bez znanego końca próby może uniemożliwić obliczenie łącznego czasu rzeczywistego. Brakująca wartość nie oznacza zerowego kosztu.

\texttt{selected\_inference\_seconds} sumuje zapisane przebiegi predykcji konfiguracji o~dodatniej wadze. Nie jest nowym pomiarem końcowej latencji całego zespołu i~nie obejmuje automatycznie wszystkich kosztów składania wyników. Czas wyboru z~pamięci podręcznej raportowany jest osobno. Koszt znalezienia kandydatów jest wspólny dla D, T i~T+E, dlatego sumowanie tych wierszy prowadziłoby do jego wielokrotnego naliczenia.

Warunki wykonania obejmują cztery wątki CPU oraz limit partii predykcji 4096, przekazywany lokalnym MLP, FT-Transformerowi i~TabNet. W~dwóch pierwszych metodach obowiązuje również ograniczenie partią uczenia. Natywne integracje i~TabPFN określają podział predykcji na partie we własnych implementacjach. Zakres pomiarów zestawiono w~tabeli~\ref{tab:zakres-pomiarow}.

\begin{table}[tbp]
\centering\small
\caption{Zakres pomiarów pamięci podczas dopasowania.}
\label{tab:zakres-pomiarow}
\begin{tabularx}{\linewidth}{P{0.18\linewidth}YY}
\toprule
Miara & Zakres & Granica interpretacji \\
\midrule
RSS & Pamięć rezydentna procesu na CPU & Nie jest pomiarem całej maszyny \\
Pamięć PyTorch & Pamięć przydzielona i~rezerwowana na GPU & Nie obejmuje wszystkich alokacji procesu ani innych użytkowników GPU \\
Dostępność śledzenia & Znacznik w~bloku \texttt{memory\_usage} & Brak pomiaru nie oznacza zerowego zużycia \\
\bottomrule
\end{tabularx}
\end{table}

Pomiary pamięci są zapisywane w~bloku \texttt{memory\_usage} rekordu konfiguracji i~obejmują dopasowanie. Nie wyznaczają szczytu całej próby ani pamięci testowej predykcji. Osobny pomiar inferencji jest szczególnie istotny dla KNN i~TabPFN, które zachowują zbiór odniesienia lub kontekst zadania. Zewnętrzne pretrenowanie TabPFN pozostaje poza kosztem lokalnego wykonania.

Współdzielenie GPU wymaga dodatkowej kontroli: zapis modelu urządzenia i~jego pamięci nie określa obciążenia przez inne procesy. Przy interpretacji czasu oraz niepowodzeń należy powiązać próby z~dostępnymi zapisami warunków wykonania. Sam komunikat o~wyczerpaniu pamięci nie pozwala przypisać całego zużycia badanemu modelowi.

Moduł metryk zapisuje nieokreślone ROC AUC jako brak wartości, lecz nie zapisuje osobnego pola z~przyczyną. Dla precyzji, czułości i~F1 stosuje natomiast konwencję \texttt{zero\_division=0}: przy zerowym mianowniku zwraca zero. Rozróżnienie tych przypadków od zerowej wartości poprawnie określonej miary wymaga uwzględnienia liczności klas i~macierzy pomyłek.

\subsection{Weryfikacja poprawności platformy}
\label{sec:kwalifikacja}

Weryfikacja platformy dotyczy trzech poziomów: kontraktów danych, poprawności przebiegu eksperymentu i~integralności zapisanych artefaktów. Kontrole wymiarów i~prawdopodobieństw wykrywają błędy interfejsów, natomiast zgodność identyfikatorów wierszy i~grup pozwala sprawdzić przepływ informacji między treningiem, walidacją i~testem.

Test syntetyczny \texttt{smoke} obejmuje niewielki przebieg z~regresją logistyczną i~dwuepokowym MLP. Służy kontroli uruchamiania, zapisu, wznowienia i~raportowania. Jego kryterium stanowi spójność artefaktów i~przebiegu, a~nie jakość predykcji. Test wznowienia wymaga zachowania ukończonych rekordów i~ponowienia wyłącznie tych konfiguracji, dla których nie istnieje kompletny, zgodny zapis.

Poprawność podziału grupowego oznacza rozłączność pacjentów na każdym poziomie, zachowanie wymaganych klas oraz dokładnie jednokrotne pokrycie obserwacji predykcjami OOF. Sam brak wspólnych indeksów wierszy nie wystarcza, ponieważ różne pobyty mogą dotyczyć tej samej osoby. Analogicznie numeryczna poprawność macierzy prawdopodobieństw nie zastępuje zgodności znaczenia jej kolumn.

Dowód weryfikacji odnosi się do konkretnej rewizji programu, danych i~środowiska. Zmiana tych składników ogranicza możliwość przenoszenia wcześniejszych sprawdzeń na nowe wykonanie. Szczególne znaczenie mają granice pamięci, asynchroniczność GPU oraz zależność kosztu inferencji TabPFN od rozmiaru kontekstu.

\subsection{Środowisko wykonania i procedura kampanii}
\label{sec:wykonanie-kampanii}

Podstawą odtworzenia środowiska jest Python~3.11 na Linux x86-64 oraz plik \path{uv.lock}. Istotne deklarowane wersje zestawiono w~tabeli~\ref{tab:wersje}. Wersje pakietów identyfikują implementacje i~nie są tożsame z~datami publikacji metod. Wersja pakietu \texttt{tabpfn} nie zastępuje identyfikacji konkretnego checkpointu modelu.

\begin{table}[tbp]
\centering\small
\caption{Wybrane przypięcia zależności środowiska. Pełna tożsamość wymaga zachowania pliku blokady i~manifestu uruchomienia.}
\label{tab:wersje}
\begin{tabularx}{\linewidth}{YY}
\toprule
Składnik & Wersja \\
\midrule
TabArena & 0.1.0 \\
scikit-learn & 1.9.0 \\
XGBoost & 3.2.0 \\
CatBoost & 1.2.10 \\

AutoGluon core i~tabular & 1.6.3b20260918 \\
PyTorch & 2.13.0 \\
\texttt{pytabkit} & 1.8.0 \\
\texttt{rtdl-revisiting-models} & 0.0.2 \\
\texttt{pytorch-tabnet} & 4.1.0 \\
\texttt{tabm} & 0.0.3 \\
\texttt{tabpfn} & 9.2.0 \\
\bottomrule
\end{tabularx}
\end{table}

Odtworzenie kampanii wymaga kolejno przygotowania środowiska z~zamrożonych zależności, audytu i~przygotowania danych, udostępnienia checkpointu oraz sprawdzenia GPU. Rewizję, nazwę pliku i~hash wag wskazuje konfiguracja. Powiązanie tych identyfikatorów z~manifestem wykonania rozróżnia modele korzystające z~tej samej wersji pakietu, lecz innych wag.

Manifesty, indeksy, historia prób i~predykcje tworzą podstawę odtworzenia raportów. Interpretacja naukowa pozostaje odrębna od działania raportera: wymaga określenia populacji porównania, znaczenia różnic i~zakresu niepewności.
